\documentclass{article}
\usepackage[T1]{fontenc}
\usepackage{lmodern}
\usepackage{graphicx}
\usepackage{amsmath,amssymb}
\usepackage[a4paper,margin=2cm]{geometry}
\usepackage[absolute,overlay]{textpos}
\setlength{\TPHorizModule}{1mm}
\setlength{\TPVertModule}{1mm}
\usepackage[indonesian]{babel}
\usepackage{hyperref}
\setlength{\emergencystretch}{2em}
% unit: Halaman 39

\begin{document}

% === Halaman 39 (python/native) ===

Ranah Steganografi

Berdasarkan  ranah operasinya, metode-metode steganografi dapat dibagi menjadi 

dua kelompok 

• Spatial (time) domain methods

    Memodifikasi langsung nilai byte dari cover-object (nilai byte merepresentasikan 

intensitas/warna pixel atau amplitudo)

 Contoh: Metode LSB

• Tranform domain methods

    Memodifikasi hasil transformasi sinyal dalam ranah transform (hasil  transformasi 

dari ranah spasial ke ranah lain (misalnya ranah frekuensi). 

 Contoh: Metode Spread Spectrum

39

\end{document}
